\item \model{DeepSeek-V4}:
پیکره داده‌های مرحله پیش‌آموزش (\en{Pre-Training}) این خانواده از مدل‌ها—شامل دو نسخه پایه \en{DeepSeek-V4-Flash} و \en{DeepSeek-V4-Pro}—با هدف ارتقای بنیادین ظرفیت‌های استدلالی، مهارت‌های کدنویسی، گسترش دانش چندزبانه و پشتیبانی بومی از طول زمینه یک میلیون توکن طراحی و گردآوری شده است \cite{deepseek_v4_2026}. این پیکره بر بستر تجارب مدل \en{DeepSeek-V3} \cite{deepseek2024deepseekv3} استوار بوده و از پالایش‌های الگوریتمی پیشرفته بهره می‌برد.

\paragraph*{کلیات و شالوده داده‌ها (\en{Foundational Data Composition})}
مجموع حجم کل توکن‌های مصرف‌شده در فاز پیش‌آموزش بیش از ۳۲ تریلیون توکن ($>32\text{T}$) است. مدل \en{DeepSeek-V4-Flash} بر روی $32\text{T}$ توکن و مدل بزرگ‌تر \en{DeepSeek-V4-Pro} بر روی $33\text{T}$ توکن آموزش دیده‌اند. این پیکره تنوع گسترده‌ای از متون وب پالایش‌شده، کدهای برنامه‌نویسی، ادبیات ریاضیات، اسناد علمی طویل و متون چندفرهنگی را در بر می‌گیرد.

\paragraph*{پالایش داده‌های وب و جلوگیری از فروپاشی مدل (\en{Web Data Filtering \& Model Collapse})}
در بخش متون استخراج‌شده از وب، فیلترهای الگوریتمی نوینی برای شناسایی و حذف متون قالبی (\en{Templated}) و محتواهای تولیدشده به صورت خودکار توسط سایر مدل‌های هوش مصنوعی اعمال شده است.
\begin{itemize}
    \item \textbf{تشریح پدیده فروپاشی مدل (\en{Model Collapse}):} مطابق یافته‌های ژو و همکاران \cite{zhu2024synthesize}، آموزش بازگشتی مدل‌ها بر داده‌های سنتتیک و تولیدشده توسط مدل‌های نسل پیشین، منجر به افت واریانس بازنمایی، تحلیل رفتن دم‌های توزیع احتمالاتی (\en{Long-Tail Distribution}) و در نهایت زوال تنوع زبانی و استدلالی مدل می‌شود. فیلترینگ سخت‌گیرانه داده‌های وب در \en{DeepSeek-V4} مانع از وقوع این پدیده شده است.
\end{itemize}

\paragraph*{پیکره تخصصی ریاضیات، کدنویسی و داده‌های عاملیتی (\en{Code, Math \& Agentic Data})}
متون ریاضیاتی و مخازن کد هسته اصلی داده‌ها را تشکیل می‌دهند. ویژگی منحصر‌به‌فرد این نسخه، تزریق \textbf{داده‌های عاملیتی (\en{Agentic Data})} در فاز میانی آموزش (\en{Mid-Training}) است.
\begin{itemize}
    \item \textbf{داده‌های عاملیتی (\en{Agentic Data}):} به توالی‌هایی اطلاق می‌شود که شامل تعاملات گام‌به‌گام با محیط، فراخوانی ابزارها (\en{Tool Calling})، تحلیل خروجی مفسر کد و بازخورد خطاها هستند. این داده‌ها مدل پایه را پیش از ورود به فاز پس‌آموزش، با منطق عامل‌های خودمختار آشنا می‌سازند.
    \item \textbf{فاز میانی آموزش (\en{Mid-Training}):} مرحله‌ای پیوسته در اواخر پیش‌آموزش عمومی که در آن توزیع داده‌ها به سمت سناریوهای پیچیده استدلالی و تعاملی سوق داده می‌شود تا فضای پنهان مدل برای یادگیری پس‌آموزش تقویت گردد.
\end{itemize}

\paragraph*{داده‌های چندزبانه و دانش دم‌بلند (\en{Multilingual \& Long-Tail Knowledge})}
به منظور بهبود بازنمایی مفاهیم در زبان‌های مختلف و ارتقای فهم ظرایف فرهنگی، سهم داده‌های چندزبانه به طور چشمگیری افزایش یافته است. این راهبرد یادگیری حقایق نادر و دانش دم‌بلند (\en{Long-Tail Knowledge})—دانشی که فراوانی تکرار اندکی در اسناد وب عمومی دارد—را بهبود بخشیده است.

\paragraph*{گردآوری متون طویل و آکادمیک (\en{Long-Document Curation})}
با توجه به هدف‌گذاری پنجره زمینه یک میلیونی ($1\text{M}$)، بخش ویژه‌ای از خط لوله داده به متون طولانی با پیوستگی معنایی بالا نظیر مقالات علمی، گزارش‌های فنی، کتاب‌ها و تحلیل‌های پژوهشی ساختاریافته اختصاص یافته است تا ساختار وابستگی‌های دوربرد در مدل پایدار شود.

\paragraph*{پیش‌پردازش، توکن‌سازی و راهبردهای قالب‌بندی (\en{Pre-processing \& Formatting})}
فرآیند آماده‌سازی توکن‌ها با تکنیک‌های زیر اجرا گردیده است:
\begin{itemize}
    \item \textbf{فضای واژگان (\en{Vocabulary}):} حفظ ساختار توکنایزر \en{DeepSeek-V3} با اندازه واژگان ۱۲۸ هزار توکن ($128\text{K}$) و معرفی توکن‌های کنترلی جدید جهت تعیین مرزهای بافتار (\en{Context Construction}).
    \item \textbf{پر کردن بخش میانی (\en{Fill-in-the-Middle - FIM}):} استفاده از سازوکار تقطیع و بازآرایی متن به قالب پیشوند، پسوند و میانه برای ایجاد قابلیت درک دوطرفه در متن و تکمیل کدها بر اساس رویکرد \cite{deepseek2024deepseekv3}.
    \item \textbf{بسته‌بندی اسناد (\en{Document Packing}):} با الهام از راهبرد دینگ و همکاران \cite{ding2024fewer}، جهت جلوگیری از هدررفت محاسباتی ناشی از پدگذاری (\en{Padding}) و جلوگیری از آسیب ناشی از برش ناگهانی انتهای متون (\en{Sample Truncation})، اسناد کوتاه متعدد درون یک پنجره زمینه واحد ادغام می‌شوند.
    \item \textbf{ماسک‌گذاری توجه در سطح نمونه (\en{Sample-Level Attention Masking}):} برخلاف رویکرد آموزش پیوسته پیشین، توجه بین‌نمونه‌ای در یک توالی ادغام‌شده به طور کامل توسط ماسک مسدود می‌شود تا از نشت بافتار (\en{Context Contamination}) بین اسناد غیرمرتبط جلوگیری گردد.
\end{itemize}

\paragraph*{برنامه زمانی طول زمینه و اندازه دسته (\en{Curriculum \& Scheduling})}
فرآیند پیش‌آموزش داده‌ها دارای یک برنامه زمانی فزاینده (\en{Curriculum}) در دو بعد طول زمینه و اندازه دسته بوده است:
\begin{itemize}
    \item \textbf{برنامه طول دنباله:} آموزش مدل‌ها ابتدا با طول دنباله $4\text{K}$ آغاز شده و به تدریج به $16\text{K}$، $64\text{K}$ و نهایتاً به $1\text{M}$ توکن ارتقا یافته است:
    \begin{equation}
        L \in \{4\text{K} \to 16\text{K} \to 64\text{K} \to 1\text{M}\}
    \end{equation}
    \item \textbf{زمان‌بندی اندازه دسته:} اندازه دسته توکن‌ها از مقادیر اندک آغاز و تا سطوح پایدار مقیاس‌بندی شده است؛ برای نسخه \en{Flash} تا سقف $75.5\text{M}$ توکن و برای نسخه \en{Pro} تا سقف $94.4\text{M}$ توکن در هر گام بهینه‌سازی.
    \item \textbf{گذار توجه متراکم به کم‌پشت:} در طول ۱ تریلیون توکن نخست ($1\text{T}$)، مدل با سازوکار توجه متراکم (\en{Dense Attention}) گرم‌سازی شده و از طول زمینه $64\text{K}$ به بعد، سازوکار توجه کم‌پشت (\en{Sparse Attention}) فعال و جایگزین شده است.
\end{itemize}

\paragraph*{تحلیل‌های آماری و فرمولاسیون ریاضی داده‌ها (\en{Statistical \& Mathematical Analysis})}
حجم توکن کل پردازش‌شده ($\mathcal{D}$) در کل مراحل بهینه‌سازی گام $t \in \{1, \dots, T\}$ از رابطه زیر پیروی می‌کند:
\begin{equation}
    \mathcal{D} = \sum_{t=1}^{T} B(t)
\end{equation}
که در آن برای مدل‌های مذکور مقادیر نهایی برابرند با:
\begin{equation}
    \mathcal{D}_{\text{Flash}} = 3.2 \times 10^{13} \text{ \en{tokens}}, \quad \mathcal{D}_{\text{Pro}} = 3.3 \times 10^{13} \text{ \en{tokens}}
\end{equation}

در فرآیند بسته‌بندی اسناد با اعمال ماسک‌گذاری سطحی، ماتریس ماسک علّی-نمونه‌ای $M \in \mathbb{R}^{L_{\text{seq}} \times L_{\text{seq}}}$ بر روی شباهت‌های لایه توجه به صورت زیر تعریف می‌شود:
\begin{equation}
    M_{i,j} = \begin{cases} 
        0 & \text{if } j \le i \text{ and } \text{DocID}(i) = \text{DocID}(j) \\ 
        -\infty & \text{otherwise} 
    \end{cases}
\end{equation}
که در آن $\text{DocID}(i)$ شناسه سندی است که توکن $i$-ام به آن تعلق دارد. در نتیجه، ماتریس توجه همواره ارتباط متقابل اسناد ادغام‌شده را به صفر مطلق می‌رساند:
\begin{equation}
    \text{Attention}(Q, K, V) = \text{Softmax}\left(\frac{QK^T}{\sqrt{d_k}} + M\right)V
\end{equation}

همچنین در زیرساخت موازی‌سازی بافتار (\en{Context Parallelism}) همراه با فشرده‌سازی توجه، به منظور پیشگیری از هدررفت توکن‌ها در مرزهای توالی به واسطه ضریب فشرده‌سازی لایه‌ها ($m=4$ در \en{CSA} و $m'=128$ در \en{HCA})، داده‌های ورودی بر مبنای ک.م.م این ضرایب هم‌تراز می‌گردند:
\begin{equation}
    L_{\text{block}} = k \cdot \text{lcm}(m, m') = k \cdot \text{lcm}(4, 128) = 128k, \quad k \in \mathbb{N}
\end{equation}
این هم‌ترازی تضمین می‌نماید که اتلاف توکن ناشی از باقیمانده طول دنباله ($l_{\text{seq}} \pmod m$) به حداقل ممکن کاهش یابد.